% =============================================================================
%  Chapter 3 — Phương pháp đề xuất
%  Section: Mở đầu chương
% =============================================================================

Chương này trình bày quá trình phân tích và thiết kế kiến trúc hệ thống GraphRAG phục vụ việc khai thác tri thức trong lĩnh vực nghệ thuật Chèo. Trên cơ sở các nền tảng lý thuyết đã trình bày ở Chương 2, chương này tập trung mô tả cách thức xây dựng đồ thị tri thức và tích hợp nó với mô hình ngôn ngữ lớn nhằm xây dựng một hệ thống hỏi đáp có khả năng khai thác tri thức có cấu trúc. Nội dung chương 3 bao gồm mô tả kiến trúc tổng thể của hệ thống và phân tích các thành phần chính như tầng tri thức, phân hệ truy xuất đồ thị (G-Retrieval) và phân hệ tạo sinh (G-Generation), qua đó làm rõ cấu trúc và cơ chế vận hành của hệ thống GraphRAG được đề xuất. Bên cạnh đó, chương cũng đề xuất một bộ dữ liệu đánh giá mới, làm cơ cơ để nghiên cứu và phân tích hiệu quả của việc áp dụng các mô hình ngôn ngữ trong khai phá miền dữ liệu truyền thống.
